\documentclass[11pt,a4paper]{article}
\usepackage[utf8]{inputenc}
\usepackage[T1]{fontenc}
\usepackage[spanish]{babel}
\usepackage{amsmath,amssymb}
\usepackage{geometry}
\geometry{margin=2.5cm}
\usepackage{siunitx}
\sisetup{output-decimal-marker={,}}
\usepackage{enumitem}
\usepackage{booktabs}
\usepackage{hyperref}

\usepackage{mathtools}
\usepackage{gensymb}

\title{\textbf{Guía de lectura de la gráfica de una onda plana}\\[6pt]
\large Material de apoyo — \textit{Diseño de Antenas}}
\author{}
\date{}

\begin{document}
\maketitle
\vspace{-2.5em}

\section{Qué representa la onda plana}
Una \textbf{onda plana uniforme} es una onda electromagnética cuyo campo
eléctrico y campo magnético son perpendiculares entre sí y perpendiculares a la
dirección de propagación. Si se propaga en $+\hat{x}$, su expresión es
\begin{equation}
  \mathbf{E}(x,t) = A\cos(\omega t - kx)\,\hat{y},
  \qquad\qquad
  \mathbf{H}(x,t) = \frac{A}{\eta}\cos(\omega t - kx)\,\hat{z},
  \label{eq:onda}
\end{equation}
donde el argumento $(\omega t - kx)$ es la \textbf{fase} de la onda. La gráfica
que acompaña este documento muestra esta función de dos variables: la misma
onda vista simultáneamente respecto de la \emph{distancia} y respecto del
\emph{tiempo}.

\subsection{Parámetros que aparecen en la figura}
\begin{center}
\begin{tabular}{@{}lll@{}}
\toprule
\textbf{Símbolo} & \textbf{Significado} & \textbf{Valor (ejemplo 300\si{\mega\hertz})}\\
\midrule
$A$          & Amplitud del campo eléctrico      & $1\ \si{\volt\per\metre}$\\
$f$          & Frecuencia de la onda             & $300\ \si{\mega\hertz}$\\
$\omega = 2\pi f$ & Frecuencia angular          & $2\pi(300\times10^{6})\ \si{\radian\per\second}$\\
$\lambda = \frac{c}{f}$ & Longitud de onda       & $\lambda = \frac{3\times10^{8}}{300\times10^{6}} = 1\ \si{\metre}$\\
$k = \frac{2\pi}{\lambda}$ & Número de onda     & $k = 2\pi\ \si{\radian\per\metre}$\\
$T = \frac{1}{f}$ & Periodo                     & $T = 3{,}33\ \si{\nano\second}$\\
$\eta = \frac{E}{H}$ & Impedancia intrínseca   & $\eta \approx 377\ \si{\ohm}$ (vacío)\\
\bottomrule
\end{tabular}
\end{center}
La velocidad de fase es la velocidad a la que viaja una cresta:
\begin{equation}
  v_{\varphi} = \frac{\omega}{k} = \lambda f = c.
  \label{eq:vel}
\end{equation}

\section{Cómo se lee el panel izquierdo (la superficie $E(x,t)$)}
El panel izquierdo es la superficie de la ecuación \eqref{eq:onda}: el eje
$x$ es \emph{distancia}, el eje $y$ es \emph{tiempo} y el eje vertical es
\emph{el valor del campo} $E$ en $\si{\volt\per\metre}$.

\begin{enumerate}[leftmargin=*, itemsep=2pt]
  \item \textbf{Como onda respecto de la distancia.} Tomar un corte a \emph{un
  tiempo fijo}, es decir, un plano paralelo al eje de distancia. Ese corte es la
  onda vista en el espacio en un instante: es la \textbf{línea roja} animada. Su
  forma es $E(x) = A\cos(kx - \omega t_0)$; es una sinusoide de periodo espacial
  $\lambda$. Distancia entre dos crestas consecutivas $\Rightarrow$ $\lambda$.
  \item \textbf{Como onda respecto del tiempo.} Tomar un corte a \emph{una
  distancia fija} $x_0$. Ese corte es el historial del campo en un punto: es la
  \textbf{línea celeste}. Su forma es $E(t) = A\cos(\omega t - kx_0)$; una
  sinusoide de periodo temporal $T$. Duración de un ciclo completo $\Rightarrow$
  $T$.
  \item \textbf{Ambas son la misma onda.} La línea roja y la celeste se cruzan:
  el punto de cruce es el valor del campo \emph{en esa distancia y ese tiempo}.
  Mientras corre la animación, la roja se desliza sobre la celeste, lo que
  muestra que un punto fijo del espacio “ve” oscilar el campo.
  \item \textbf{La cresta blanca ($x = ct$).} Une los máximos (los “arbolitos”
  verdes de la superficie, que son sencillamente los picos donde $E = +A$). Su
  diagonal indica la velocidad de propagación \eqref{eq:vel}: la onda avanza
  $1\ \si{\metre}$ cada $3{,}33\ \si{\nano\second}$.
  \item \textbf{Marcas en la figura.} $\lambda$ y $T$ quedan señaladas entre
  crestas; el texto ``valor negativo: $-A$'' muestra el semiciclo negativo (el
  campo alterna signo, propiedad esencial de la onda).
\end{enumerate}

\section{Cómo se lee el panel derecho (orientación de los campos)}
Aquí se dibujan los dos campos de la ecuación \eqref{eq:onda} por separado y
perpendiculares entre sí:
\begin{itemize}[leftmargin=*, itemsep=2pt]
  \item la \textbf{onda roja} es el campo eléctrico $E_y$, que oscila a lo
  largo del eje $y$ (arriba/abajo);
  \item la \textbf{onda azul} es el campo magnético $H_z$, que oscila a lo
  largo del eje $z$ (dentro/fuera de la pantalla);
  \item ambas avanzan hacia la derecha, en $+x$ (flecha blanca), y están
  \textbf{en fase}: máximos y ceros coinciden en el espacio.
\end{itemize}

\subsection{El relleno y la perspectiva}
El \textbf{relleno bajo la curva} (área entre la onda y su eje) es exactamente
el recurso de la integral definida: visualmente lee la magnitud de la onda en
cada punto sin esfuerzo. Cuando la animación corre, el relleno crece y decrece
con el campo: su “altura” en cualquier $x$ es el valor instantáneo de $E$ o de
$H$. Es convencional que la zona bajo $E$ quede teñida del color del campo para
diferenciar las dos ondas sin confundirse.

\subsection{La relación de amplitudes}
La impedancia intrínseca relaciona ambas amplitudes:
\begin{equation}
  \eta = \frac{E}{H} = \sqrt{\frac{\mu_0}{\varepsilon_0}} \approx 377\ \si{\ohm}
  \quad\Longrightarrow\quad
  H_m = \frac{A}{\eta}.
  \label{eq:eta}
\end{equation}
En la figura, $H$ se dibuja \emph{amplificada por $\eta$} para que su cresta
coincida con la de $E$ (a escala real el campo magnético mediría
$A/\eta \approx 2{,}65\ \si{\milli\ampere\per\metre}$, invisible frente a
$1\ \si{\volt\per\metre}$). La ecuación \eqref{eq:eta} es la que da la magnitud
real de $H$ a partir de la de $E$.

\section{Resumen de verificación}
\textbf{Para saber si se está leyendo bien la gráfica, comprobar:}
\begin{enumerate}[leftmargin=*, itemsep=2pt]
  \item \textbf{$\lambda$}: distancia horizontal entre dos picos de la onda
  roja del panel derecho (y entre dos crestas de la superficie). A
  $300\ \si{\mega\hertz}$ debe ser $1\ \si{\metre}$.
  \item \textbf{$T$}: separación temporal entre dos máximos de la línea celeste
  del panel izquierdo. Debe ser $3{,}33\ \si{\nano\second}$.
  \item \textbf{Ortogonalidad}: $E_y \perp H_z \perp \hat{x}$; en el panel
  derecho se ve que las oscilaciones usan dos ejes perpendiculares.
  \item \textbf{Fase}: en el vacío ambos campos alcanzan el máximo y el cero
  al mismo tiempo (no hay desfase).
  \item \textbf{Sentido}: un observador fijo ve el campo oscilar; un
  ``fotógrafo'' mueve el corte rojo y ve la onda viajar en $+x$, nunca en $-x$.
\end{enumerate}

\vspace{1em}
\noindent\emph{Nota:} al rotar los paneles 3D con el ratón se confirma cada
afirmación anterior desde cualquier ángulo; el relleno y las guías de $\pm A$
aseguran que la lectura de la magnitud sea siempre posible.

\end{document}